\section{A common normal-product coordinate}\label{sec:normal-product}

At a simple saddle, a slow-root collision and the lower endpoint, we
start from a selected invariant graph with a fixed normal spectral gap.
The construction below separates the base and normal equations while
preserving that graph. It gives a common tube, uniform finite phase
derivatives, and the small coupling factor needed when passing to
physical sections. Each application first constructs its selected graph.

The persistence result of
\cite[Theorem~3, p.~215]{Fenichel1971} assumes a $C^r$ field,
a compact connected properly embedded $C^r$ invariant manifold,
a hyperbolic splitting, and the source rate inequalities
$\sigma^\pm(m)<1/r$ for every point of that manifold.
We apply it to compact boundaryless extensions after establishing a
fixed normal gap and finite-order rate domination. Its direct conclusion
is invariant-manifold existence. Selection by the tail or incoming
separatrix, the sharp forcing factor, uniform phase jets, and physical
passage are adaptations proved in Sections~\ref{sec:simple-contact-tools},
\ref{sec:positive-constant}, \ref{sec:two-endpoints}, and
\ref{sec:lower-complement}.

\begin{lemma}\label{lem:uniform-normal-product}
Fix an integer $j\ge4$ and nested compact intervals
$I_0\Subset I_1\Subset I_2$. Suppose a chosen invariant graph has been
translated to $u=0$ and the equations are
\begin{equation}\label{eq:uniform-normal-input}
 \dot x=\tau v(x)+\rho c(x,u)u,\qquad
 \dot u=B(x,u)u,\qquad 0<\tau\le\tau_0,
 \quad 0\le\rho\le1.
\end{equation}
On a fixed tube over $I_2$, suppose $v,c,B$ have common
$C^{2j+6}$ bounds in $(x,u)$ and $B$ has one sign with $|B|\ge a_*>0$.
Then, after decreasing $\tau_0$ and the tube radius, there are exact
$C^j$ coordinates on a common tube over $I_0$ in which
\begin{equation}\label{eq:uniform-normal-product}
 \dot\theta=\tau v(\theta),\qquad
 \dot\eta=b(\theta)\eta,
 \qquad b(\theta)=B(\theta,0).
\end{equation}
The coordinate maps and their inverses have common bounds through order
$j$. They preserve the graph and satisfy
\[
 x=X(\theta,\eta)=\theta+\rho\eta L_1(\theta,\eta),
 \qquad u=\eta+\eta^2L_2(\theta,\eta),
\]
where $L_1,L_2$ have common $C^{j-1}$ bounds.
In particular $X_\eta=O_{C^{j-2}}(\rho)$ and the normal coordinate is
uniformly transverse to the graph. All bounds are for phase derivatives
with the field parameters fixed. Zeros of $v$, including multiple zeros,
are allowed.
\end{lemma}

\begin{proof}
 Put $\bar\tau:=|\tau|$. If $B\ge a_*$, first apply the argument to the
 reversed vector field, replacing $v$ by $-v$ so that its base scale is the
 positive $\bar\tau$. All estimates below use this $\bar\tau$; reversing the
 orientation at the end restores the stated equation with the original
 positive $\tau$. We may therefore assume $B\le-a_*$. Extend the already
 translated equations outside $I_2$ with fixed smooth cutoffs, equal to the
 original coefficients on $I_2$, keeping $u=0$ invariant, $B\le-a_*$, and
 the factor $\tau$ in its base flow. Choose the extension of $B$ by a smooth
 interpolation to a negative constant, and taper $v,c$ smoothly only outside
 $I_2$; no extension of a prescribed boundary condition is being asserted.
 This gives bounded coefficients on the whole base line. A smaller normal
 tube is forward invariant. Write $\varphi_t$ for the complete flow of
 $\dot x=\tau v(x)$ and $(x_t,u_t)$ for the full flow.
The asymptotic phase is the convergent limit
\begin{equation}\label{eq:uniform-phase-limit}
 \theta(x,u)=\lim_{t\to\infty}\varphi_{-t}(x_t)
 =x+\rho\int_0^\infty
 D\varphi_{-t}(x_t)c(x_t,u_t)u_t\,dt.
\end{equation}
Indeed $|u_t|\le |u|e^{-a_*t}$ and
$\|D^k\varphi_{\pm t}\|\le K_k(1+t)^{N_k}e^{K_k\bar\tau t}$.
The required differentiated estimates follow directly from the flow
variations. At first order they have coefficient matrix
\[
 \begin{pmatrix}
 \tau v_x+\rho c_xu&\rho(c+uc_u)\\
 B_xu&B+uB_u
 \end{pmatrix}.
\]
Variation of constants and $\int_0^\infty|u_t|dt\le |u|/a_*$
bound the longitudinal variations by a polynomial times
$e^{K\bar\tau t}$.
Every forcing term in a normal variation contains either $u_t$ or a
lower normal variation. Induction through order $j+3$ consequently
bounds the differentiated integrand in \eqref{eq:uniform-phase-limit}
by
\[
 K_j(1+t)^{N_j}
       e^{-(a_*-K_jd-K_j\bar\tau)t},
\]
where $d$ is the normal tube radius. Choose $K_jd<a_*/4$ and then
$K_j\bar\tau_0<a_*/4$. The integral and these derivatives converge
uniformly. In particular $\theta=x$ on $u=0$,
$\theta-x=\rho u L(x,u)$ with bounded finite derivatives, and
$\theta_x=1+O(d)$. Thus $(\theta,u)$ are uniformly invertible
coordinates. The flow identity in the limit gives
$\dot\theta=\tau v(\theta)$ exactly.

In these coordinates write the second equation as
$\dot u=b(\theta)u+q(\theta,u)u^2$ and put
\begin{equation}\label{eq:uniform-amplitude-limit}
 \eta=\lim_{t\to\infty}u_t
       \exp\left(-\int_0^t b(\varphi_s\theta)ds\right)
 =u\exp\left(\int_0^\infty
                    q(\varphi_t\theta,u_t)u_t\,dt\right).
\end{equation}
The same induction, now with the base equation independent of $u$,
gives an integrable majorant for the required derivatives. The last
expression is $u+O(u^2)$ and has positive normal derivative on the
smaller tube. Shifting the limit in time proves
$\dot\eta=b(\theta)\eta$. The inverse function theorem gives the
asserted common coordinate bounds. The input order $2j+6$ leaves more
than the derivatives lost in division by $u$, these two coordinate
changes, and the section changes below. Reversing time back proves
the assertion for either normal sign.
\end{proof}

A product of order $J$ requires $C^{2J+6}$ coefficients and, when a
coefficient contains one derivative of the selected graph, a
$C^{2J+7}$ graph. Fix these orders before shrinking the parameter window.
The simple-contact and positive-collision applications take $J=j+2$,
with graph order at least $2j+11$ and $C^{2j+10}$ coefficients;
this retains the $C^j$ small shear needed for physical sections and the
full collision output $N$. The lower-endpoint product and the lower
$B$ output take $J=j+1$, with graph order $2j+9$ and
$C^{2j+8}$ coefficients. Only finite phase derivatives with the field
parameters fixed are used.

To use the product for physical transitions, we need two section changes. If
$\rho=\tau$ and $|v(\theta_b)|\ge v_*>0$ at an exterior base, the
change from $\theta=\theta_b$ to $x=\theta_b$ has uniformly bounded
finite derivatives. In fact its hitting equation, divided by $\tau$,
is
\[
 \int_0^T v(\varphi_t\theta_b)dt
       +\eta_T L_1(\varphi_T\theta_b,\eta_T)=0,
 \qquad \eta_T=\eta\exp\int_0^T b(\varphi_t\theta_b)dt.
\]
Its derivative in $T$ is bounded away from zero for small $\eta$;
therefore $T=O(\eta)$ with common derivative bounds. The required
lower bound on $|v|$ is at the exterior base; $v$ may vanish inside the
collision strip.

For $\rho=1$, or after a further uniformly smooth change to physical
coordinates, write $z=X(\theta,\eta)$ and $y=H(\theta,\eta)$.
The scaled variable $\eta=\tau V$ gives uniform section bounds in $V$.
Assume again
$|v(\theta_b)|\ge v_*>0$. Suppose $X_\theta$ is a
positive unit and the normal direction is transverse to the graph,
with $a_l:=H_\eta-h_zX_\eta$ a positive unit at
$(\theta_b,0)$, where $y=h(z)$ is the physical graph. Divide
$X(\varphi_T\theta_b,\eta_T)-X(\theta_b,0)=0$ by $\tau$
after this substitution. The implicit function theorem gives bounded
derivatives in $V$ for $|V|\le d_1$, where $d_1>0$ is fixed and small.
To retain the normal factor, write
$H(\theta,\eta)-h(X(\theta,\eta))=\eta a(\theta,\eta)$ and set
$E_T(V)=\exp\int_0^{T(V)}b(\varphi_t\theta_b)dt$. Exactly,
\[
 \frac{\Phi(\tau V)-h(z_b)}{\tau}
 =V E_T(V)a\bigl(\varphi_{T(V)}\theta_b,\tau V E_T(V)\bigr).
\]
The right side has a common $C^{j-1}$ bound and derivative $a_l$ at zero.
Taylor's formula therefore gives
\begin{equation}\label{eq:uniform-physical-section}
 \Phi(\eta)=h(z_b)+a_l\eta+
       \frac{\eta^2}{\tau}E(\eta/\tau),\qquad
 \|E\|_{C^{j-3}}\le K,
 \quad \frac{|\Phi^{(k)}|}{\Phi'}\le K_k\tau^{1-k}
 \quad(2\le k\le j-1).
\end{equation}
Here $z_b=X(\theta_b,0)$ and $\Phi'$ is a positive unit after
decreasing $d_1$. Thus unscaled section derivatives need not remain
bounded as $\tau\to0$.

Finally, if $v$ has a simple zero $\theta_s$ and
$|v'(\theta_s)|$ is a positive unit, division of time by $|b(\theta)|$
gives $\dot\theta=\tau g(\theta)$ and
$\dot\eta=\operatorname{sgn}(b)\eta$. With
$\beta=g'(\theta_s)$, the coordinate
\[
 \xi=\pm(\theta-\theta_s)
 \exp\int_{\theta_s}^{\theta}
       \left(\frac{\beta}{g(t)}-\frac1{t-\theta_s}\right)dt
\]
with the sign chosen to make the incoming coordinate positive,
has common finite derivatives and satisfies
$\dot\xi=\tau\beta\xi$. The integrand has a removable singularity;
the small factor $\tau$ never enters its denominator. When the normal
direction is expanding and $\beta<0$, the passage from
$\xi=\xi_0$ to $\eta=\eta_0$ is exactly
$D(s)=\xi_0(s/\eta_0)^r$, with $r=-\tau\beta$.
This last scalar linearization is used only at simple roots.
